% Generated draft: requires mathematical and source review.
\documentclass[11pt]{article}
\usepackage[margin=25mm]{geometry}
\usepackage{fontspec}
\setmainfont{Latin Modern Roman}
\usepackage{amsmath,amssymb,mathtools,unicode-math}
\setmathfont{Latin Modern Math}
\usepackage{enumitem,needspace,xurl,hyperref}
\setlength{\parindent}{0pt}
\setlength{\parskip}{5pt}
\setlength{\emergencystretch}{3em}
\newcommand{\N}{\mathbb{N}}
\newcommand{\Z}{\mathbb{Z}}
\newcommand{\Q}{\mathbb{Q}}
\newcommand{\R}{\mathbb{R}}
\newcommand{\C}{\mathbb{C}}
\newcommand{\F}{\mathbb{F}}
\newcommand{\A}{\mathbb{A}}
\newcommand{\PP}{\mathbb{P}}
\newcommand{\E}{\mathbb{E}}
\newcommand{\eps}{\varepsilon}
\newcommand{\dd}{\,\mathrm{d}}
\begin{document}
\section*{UnsolvedMath ID: 30003201}
\textbf{Source problem code:} catalogue entry
\subsection{Definitions and mathematical statement}
Work with complex \(C^*\)-algebras. A \(C^*\)-algebra is separable if it has a countable norm-dense subset, unital if it has a multiplicative identity, and simple if it is nonzero and has no nonzero proper closed two-sided ideals. Interpret amenability in its usual \(C^*\)-algebraic sense, equivalently nuclearity. Let \(\mathcal Z\) denote the Jiang--Su algebra; \(\mathcal Z\)-stability means \(A\cong A\otimes_{\min}\mathcal Z\) as \(C^*\)-algebras. The UCT condition means that \(A\) satisfies the universal coefficient theorem for \(KK\)-theory.

A UHF algebra is a unital inductive limit of full matrix algebras with unital connecting maps. It is of infinite type if its supernatural number \(\prod_p p^{n_p}\) has \(n_p\in\{0,\infty\}\) for every prime \(p\), with at least one exponent equal to \(\infty\). Write \(\operatorname{gTR}(B)\) for the generalized tracial rank of \(B\). The supplied record names this technical invariant but does not provide its defining approximation class or detailed convention; these must be checked against the cited report before a fully self-contained definition can be fixed.

Using that source-dependent generalized-tracial-rank predicate, define
\[
\mathcal N_1=\left\{A\;\middle|\;
\begin{gathered}
A\text{ is a separable, unital, simple, amenable }C^*\text{-algebra},\\
A\cong A\otimes_{\min}\mathcal Z,\qquad A\text{ satisfies the UCT},\\
\exists\,U\text{ an infinite-type UHF algebra such that}
\quad\operatorname{gTR}(A\otimes_{\min}U)\leq 1
\end{gathered}
\right\}.
\]
The existentially quantified algebra \(U\) may depend on \(A\). The bound is at most one, so the definition does not require generalized tracial rank to equal one.

\emph{Problem.} Determine the scope of \(\mathcal N_1\) and characterize which \(C^*\)-algebras belong to it. The supplied formulation asks both how large this class is and which algebras it contains, but specifies neither a notion of size nor a particular proposed membership criterion. Thus it is a characterization problem rather than a single stated conjectural equivalence. The UCT hypothesis is included in the supplied corrected statement, although it is absent from the supplied original wording.
\subsection{Short English statement}
Describe the class \(\mathcal N_1\) of separable, unital, simple, amenable, \(\mathcal Z\)-stable \(C^*\)-algebras satisfying the UCT whose generalized tracial rank becomes at most one after tensoring with some infinite-type UHF algebra. Determine its extent and membership criteria. The intended meaning of size and the detailed generalized-tracial-rank convention require clarification from the source.
\subsection{Sources}
\begin{itemize}
\item Supplied local UnsolvedMath registry record, permanent ID \(30003201\), title: Scope of the Generalized Tracial-Rank-One Class; source problem code \texttt{catalogue entry}; collection: Oberwolfach Reports Open Problems. No individual proposer is supplied; the proposed year is recorded as 2016.
\item Supplied report citation, with mathematical title markup normalized: \(C^*\)-Algebras (2017), \emph{Oberwolfach Reports}, pp. 2269--2345, DOI \texttt{10.4171/owr/2016/40}. \url{https://doi.org/10.4171/owr/2016/40}. The displayed year 2017 and the 2016 year in the DOI and proposed-year metadata are retained without reconciliation.
\item Supplied literature source: OWR Operator Algebras report. \url{https://ems.press/content/serial-article-files/46645?nt=1}.
\item Supplied literature source: Gardella and Perera Cuntz semigroups survey. \url{https://doi.org/10.4171/EMSS/84}.
\item Literature assessment dated exactly \texttt{2026-08-21}: the supplied curation labels the problem partially solved and reports strong classification and realization theorems for \(\mathcal N_1\), while stating that a complete intrinsic description of its members is not known. This is supplied dated curation derived from the OWR material, not an independently verified literature assessment.
\item The registry labels the statement \texttt{corrected\_verified} and reports that the official report was checked to restore the UCT hypothesis, tensor factor, generalized-tracial-rank condition, and both scope questions. This verification claim belongs to the supplied record; the cited sources have not been independently consulted for this draft.
\end{itemize}
\end{document}
